% !TEX program = xelatex
% ============================================================
%  عرض مؤتمر — Conference Talk Template
%  قالب عرض تقديمي لمحاضرة في مؤتمر علمي/أكاديمي
%  Compile with: xelatex main.tex (run twice)
% ============================================================
%  Features:
%    - 16:9 widescreen, clean default beamer theme with accent color
%    - Arabic RTL via polyglossia + Amiri font
%    - English terms wrapped with \textEN{...}
%    - Title, outline, motivation, background, methodology, results,
%      discussion, conclusions, references, Q&A slides
%    - Table and figure placeholders with TODO comments
% ============================================================

\documentclass[aspectratio=169,11pt]{beamer}

% ---- Arabic / RTL language support ----
\usepackage{polyglossia}
\setdefaultlanguage[calendar=gregorian,hijricorrection=1]{arabic}
\setotherlanguage[variant=british]{english}

% ---- Fonts (beamer uses sans by default, so set sans to Amiri too) ----
\usepackage[no-math]{fontspec}
\setmainfont[Scale=1]{Amiri-Regular.ttf}
\setsansfont[Scale=1]{Amiri-Regular.ttf}
\newfontfamily\arabicfont[Script=Arabic,Scale=0.95]{Amiri-Regular.ttf}
\newfontfamily\arabicfontsf[Script=Arabic,Scale=0.95]{Amiri-Regular.ttf}
\newfontfamily\englishfont[Scale=0.95]{TeX Gyre Termes}
\let\textEN\textenglish

% ---- Core packages ----
\usepackage{graphicx}
\usepackage{booktabs}
\usepackage{tabularx}
\usepackage{tikz}
\usepackage{amsmath}
\usepackage{amssymb}

% ---- Theme & appearance ----
\usetheme{default}
\setbeamertemplate{navigation symbols}{}
\setbeamersize{text margin left=10mm, text margin right=10mm}

% ---- Accent color ----
\definecolor{accent}{HTML}{2E5C8A}
\setbeamercolor{frametitle}{fg=accent}
\setbeamercolor{title}{fg=accent}
\setbeamercolor{alerted text}{fg=accent}
\setbeamercolor{structure}{fg=accent}
\setbeamertemplate{frametitle}[default][center]

% ---- Metadata ----
% TODO: استبدل العناوين والأسماء بالقيم الفعلية
\title{عنوان العرض}
\subtitle{عنوان فرعي للموضوع}
\author{اسم المتحدث}
\institute{الجامعة}
\date{اسم المؤتمر، شهر سنة}

\begin{document}

% ============================================================
% TITLE SLIDE
% ============================================================
\begin{frame}[plain]
\titlepage
\end{frame}

% ============================================================
% OUTLINE
% ============================================================
\begin{frame}{المحتويات}
\tableofcontents
\end{frame}

% ============================================================
% MOTIVATION
% ============================================================
\section{الدوافع}
\begin{frame}{الدوافع}
% TODO: اذكر المشكلة وأهميتها
\begin{itemize}
  \item المشكلة البحثية وما الذي يحفز دراستها.
  \item الفجوة المعرفية الحالية في المجال \textEN{(research gap)}.
  \item الأهمية العملية والنظرية للعمل.
  \item الأسئلة البحثية الرئيسية.
\end{itemize}
\end{frame}

% ============================================================
% BACKGROUND
% ============================================================
\section{الخلفية النظرية}
\begin{frame}{الخلفية النظرية}
% TODO: اعرض الأعمال السابقة باختصار
\begin{itemize}
  \item المراجعة الموجزة للأعمال السابقة \textEN{(related work)}.
  \item المفاهيم الأساسية التي يعتمد عليها العرض.
  \item القيود في الأساليب الحالية.
\end{itemize}
\end{frame}

% ============================================================
% METHODOLOGY
% ============================================================
\section{المنهجية}
\begin{frame}{المنهجية}
% TODO: صف طريقة العمل
\begin{itemize}
  \item تصميم الدراسة \textEN{(study design)}.
  \item مصادر البيانات وأدوات الجمع.
  \item طرق التحليل المستخدمة.
\end{itemize}
\end{frame}

% ============================================================
% RESULTS
% ============================================================
\section{النتائج}
\begin{frame}{النتائج — نظرة عامة}
% TODO: اعرض النتائج الرئيسية
\begin{itemize}
  \item النتيجة الأولى مع الأهمية.
  \item النتيجة الثانية مقارنة بالمتوقع.
  \item ملاحظات غير متوقعة \textEN{(unexpected findings)}.
\end{itemize}
\end{frame}

\begin{frame}{النتائج — جدول}
% TODO: استبدل الأرقام بالقيم الفعلية
\begin{table}[h]
\centering
\begin{tabularx}{0.85\textwidth}{lXX}
\toprule
المجموعة & المتوسط & الانحراف المعياري \\
\midrule
المجموعة أ & 12.4 & 2.1 \\
المجموعة ب & 18.7 & 1.9 \\
المجموعة ج & 15.2 & 3.0 \\
\bottomrule
\end{tabularx}
\end{table}
\end{frame}

\begin{frame}{النتائج — شكل توضيحي}
% TODO: استبدل بمسار ملف الشكل الفعلي
\begin{center}
% \includegraphics[width=0.7\textwidth]{figure.png}
\begin{tikzpicture}
\draw[accent, thick] (0,0) rectangle (8,4);
\node at (4,2) {\large [موضع الشكل / Figure placeholder]};
\end{tikzpicture}
\end{center}
\end{frame}

% ============================================================
% DISCUSSION
% ============================================================
\section{المناقشة}
\begin{frame}{المناقشة}
% TODO: ناقش التفسيرات والقيود
\begin{itemize}
  \item تفسير النتائج في ضوء الفرضيات.
  \item القيود المنهجية \textEN{(limitations)}.
  \item التعميم والمجالات المستقبلية.
\end{itemize}
\end{frame}

% ============================================================
% CONCLUSIONS
% ============================================================
\section{الاستنتاجات}
\begin{frame}{الاستنتاجات}
% TODO: لخّص المساهمات الرئيسية
\begin{itemize}
  \item المساهمة الأولى للعمل.
  \item المساهمة الثانية للعمل.
  \item التوصيات للدراسات المستقبلية.
\end{itemize}
\end{frame}

% ============================================================
% REFERENCES
% ============================================================
\section{المراجع}
\begin{frame}{المراجع}
% TODO: أضف المراجع الفعلية
\begin{itemize}
  \item المؤلف، \textEN{عنوان الورقة}، المؤتمر، السنة.
  \item المؤلف، \textEN{عنوان الورقة}، المجلة، السنة.
  \item المؤلف، \textEN{عنوان الكتاب}، الناشر، السنة.
\end{itemize}
\end{frame}

% ============================================================
% THANK YOU / Q&A
% ============================================================
\begin{frame}[plain]
\begin{center}
\vspace{1.5cm}
{\Huge\bfseries\color{accent} شكراً لكم}\\[1cm]
{\Large الأسئلة والنقاش}\\[0.8cm]
{\large اسم المتحدث}\\
{\normalsize البريد الإلكتروني: \textEN{email@example.com}}
\end{center}
\end{frame}

\end{document}
